% Self-contained section: the recipe of the AI sectors by data-center generation.
% Compile standalone with dc_recipe_engineering.tex, or \input into a document that loads
% amsmath, booktabs, array, longtable, graphicx and hyperref.

\section{The recipe of the AI sectors}\label{sec:recipe}
This note describes what it takes to build and run an AI data center. It gives a cost
recipe for five generations of facilities, from pre-2015 enterprise halls to 2027+ designs. For
each, it covers the facility, the IT equipment and the running costs, down to 6-digit NAICS
suppliers, so the same recipe can feed the 37-sector model and the 2017 detailed benchmark
tables. It prefers specialised sources (cost consultants, engineering white papers, national
labs, standards bodies, company filings, measurement studies) and says where independent sources
agree.

Seven results change the recipe. They are implemented in \texttt{ai\_recipe.py}, in the project's data-construction folder (baseline A3). The 2026 Bank of Italy rack recipe is kept as \texttt{ai\_recipe\_cnt2026.py}:
\begin{enumerate}
  \item The facility investment column needs installation labour. It becomes roughly
    Construction 49, Electrical 26, Machinery 19, professional services 6.5, not 33/40/27.
  \item Facility running costs other than electricity are mostly repairs, taxes, software,
    security and insurance. Own labour is 10--16 percent of them, not 50.
  \item Water is negligible in new designs.
  \item Facility depreciation is 0.078 (0.085--0.09 if plant depreciates at 0.10), not 0.075.
  \item IT capital must include networking and storage, 20--25 percent on top of the rack.
  \item IT depreciation of about 0.26 is supported. It is driven by obsolescence, not failure.
  \item The wholesale margin on IT purchases is about 6 percent, not 22.
\end{enumerate}
Newer facilities trade water and labour for capital. They cost more per MW, use almost no water,
and need a tenth of the staff of legacy halls.


\subsection{Scope, sources and generations}

\paragraph{Objects.} Two sectors share out the cost of AI services:
\begin{itemize}
  \item \textbf{AI\_facilities} ($f$) owns the shell and the electrical and mechanical plant,
    buys electricity and runs the building, and sells powered space.
  \item \textbf{AI\_compute} ($z$) owns the servers, networking and storage, buys space, and
    sells AI services.
\end{itemize}
A recipe is each sector's purchases by supplying industry, its labour, its capital by asset and
service life, and its investment column (which industries produce its capital goods). It is
expressed per MW of IT load; Section~\ref{sec:perdollar} converts it to per \$1.

\paragraph{Source hierarchy.} In order of preference:
\begin{enumerate}
  \item Cost consultants, engineering and vendor white papers, national labs and standards
    bodies, specialist analysts, and measurement studies.
  \item Company filings.
  \item Economics papers, only when they are primary.
\end{enumerate}
Each observation in the CSV files records whether it is primary or a re-quote, and traces
re-quotes to their origin. Anything not checked at source is marked UNVERIFIED there.

\paragraph{Generations.} Table~\ref{tab:arch} gives the five generations, with specifications
corrected against the sources.

\begin{table}[ht]\centering\small
\caption{Data-center generations}\label{tab:arch}
\begin{tabular}{>{\raggedright\arraybackslash}p{1.0cm}>{\raggedright\arraybackslash}p{2.4cm}>{\raggedright\arraybackslash}p{2.0cm}>{\raggedright\arraybackslash}p{3.2cm}>{\raggedright\arraybackslash}p{4.8cm}}
\toprule
Code & Generation & kW per rack & Cooling & IT \\
\midrule
A0 & Legacy enterprise and retail colocation, pre-2015 & 4--8 & Chillers, CRAH units, cooling towers & 2-socket CPU servers \\
A1 & Hyperscale cloud, air-cooled, 2015--22 & 8--15 & Outside air, evaporative or adiabatic & CPU servers, storage servers \\
A2 & AI, air-cooled or hybrid, 2023--24 & $\approx$40 & Rear-door heat exchangers & HGX H100 8-GPU servers \\
A3 & AI, liquid-cooled, 2025--26 & 120--145 & Direct-to-chip, closed loop; dry coolers or chillers & GB200 / GB300 NVL72 \\
A4 & Next generation, 2027+ & 330 (Rubin NVL72) to 600 (Kyber) & Liquid; 800 V DC distribution & Vera Rubin \\
\bottomrule
\end{tabular}
\end{table}

\subsection{What changes relative to the 2026 recipe}

\begin{longtable}{>{\raggedright\arraybackslash}p{2.8cm}>{\raggedright\arraybackslash}p{3.0cm}>{\raggedright\arraybackslash}p{4.8cm}>{\raggedright\arraybackslash}p{3.8cm}}
\caption{Current recipe against the evidence}\label{tab:changes}\\
\toprule
Item & 2026 rack recipe & Evidence & Main sources \\
\midrule\endfirsthead
\toprule
Item & Now & Evidence & Main sources \\
\midrule\endhead
IT vs facility share of capex & 67/33 (one rack sheet) & 68--72/28--32. Confirmed; rises to $\approx$72 once networking is included & Van Nieuwerburgh; Epoch; Fed FEDS 2025-109; Microsoft filings; SemiAnalysis \\
Facility cost per MW (A3) & Implicit in the rack sheet & \$12m hard (10.5--14), \$15m all-in (12--18). US mean rose from \$8.9m (2021) to \$11.1m (2025) & Turner \& Townsend 2019--25; Digital Realty; KPMG 2026; JLL \\
Facility investment column & Construction 33, Electrical 40, Machinery 27 & Construction 49, Electrical 26, Machinery 19, Prof.\ services 6.5 (A3). The rack sheet omits installation labour and fees & T\&T element shares; KPMG owner-furnished split; Schneider WP61, WP218 \\
Facility $\delta$ & 0.075 & 0.078 (plant 0.088, shell 0.0247). Plant lives 18--20 years; shell about 15--20\% of cost & ASHRAE; BEA; IRS; filings \\
Facility running costs other than electricity & 50/25/25 labour, Admin, Prof.\ services; unsourced & About 0.6 per \$1 of electricity. Repairs 0.19, taxes 0.13, own labour 0.10, software/IT 0.08, Admin 0.04, insurance 0.02 & KPMG; Epoch; Digital Realty; Virginia JLARC \\
Water & Ignored & Near zero in A3 with dry coolers (0.05 L/kWh). About 3\% of electricity cost in legacy halls. Price about \$6 per 1{,}000 gallons & LBNL 2024; The Green Grid; Mytton 2021; operator reports \\
Electricity & \$105k per rack-year & \$640k per MW-IT-year for A3 ($\approx$\$85k per 132 kW rack): PUE 1.15, utilisation 0.75, \$85/MWh. About 25\% of facility cost & Operator PUEs; EPRI; LBNL; EIA \\
IT capital & Rack only & Rack $\times$1.25: scale-out network 14\%, storage 4\%, integration 3\% of IT capex & SemiAnalysis; Cottier et al.; Epoch \\
IT $\delta$ & 0.25 & 0.26 (0.20--0.32): GPU servers 0.27, networking 0.20, storage 0.30. Failure adds $<$1 point & 31 SEC filings; rent indices; resale; Epoch; Meta and NCSA reliability \\
IT investment margins & Wholesale 22\% & Wholesale 6\%, transport 1\% (direct purchase from manufacturers) & BEA PEQ; distributor margins; operator practice \\
Compute's other purchases & None & 0.06--0.08 per \$ of GPU-cloud revenue: software 0.02, telecom 0.01, maintenance 0.015, services 0.025 & CoreWeave, Nebius, IREN filings; Census SAS 518 \\
\bottomrule
\end{longtable}

\subsection{Facility capital}

\subsubsection{Cost per MW and element shares}

\begin{table}[ht]\centering\small
\caption{Facility capex per MW of IT load, US, \$m}\label{tab:capex}
\resizebox{\textwidth}{!}{\begin{tabular}{lcccccc}
\toprule
 & A0 & A1 (2019/22) & A2 & A3 & A4 (projection) \\
\midrule
Hard cost (T\&T scope) & 9 (6--13) & 8.5 / 9.8 & 10.6 (9--13) & 12 (10.5--14) & 13--16 \\
All-in (adds land, site, substation, fees, contingency) & 11 (8--30) & 10 / 11.5 & 12 (10--14) & 15 (12--18) & 16--22 \\
\addlinespace
\multicolumn{6}{l}{\textit{Share of hard cost}} \\
Electrical, incl.\ equipment & 45\% & 54\% & 52\% & 48\% & 45--55\% \\
Mechanical, incl.\ equipment & 25\% & 22\% & 25\% & 33\% & 30--35\% \\
Shell, core, architectural & 20\% & 14\% & 12\% & 9\% & 7--9\% \\
General contractor preliminaries and fees & 10\% & 10\% & 10\% & 10\% & 8--10\% \\
\bottomrule
\end{tabular}}
\end{table}

Outside the hard cost, as shares of the A3 all-in cost:
\begin{itemize}
  \item land 3--5 percent (Cushman \& Wakefield: \$584k/MW);
  \item site works 3--6 percent;
  \item the customer substation 1--3 percent;
  \item professional fees for design, engineering, project management and commissioning 5--8
    percent;
  \item contingency 3--8 percent.
\end{itemize}
The liquid-cooling premium is 7--10 percent (T\&T; JLL 10). Schneider finds about zero at equal
rack density. Build time is about 18--20 months (JLL; Fed/Dodge).

\subsubsection{The investment column at detailed codes}\label{sec:faccol}

\paragraph{Equipment versus installation.} The rack-based column puts all electrical and cooling
spending on manufactured equipment. The engineering sources split it:
\begin{itemize}
  \item Electrical: 50--60 percent equipment, 25--35 installation, 15--20 design (Schneider
    WP61). Materials are 71 percent of electrical cost in WP218. KPMG's split of
    owner-furnished equipment agrees.
  \item Mechanical: about 65 percent equipment.
\end{itemize}
Table~\ref{tab:faccol} builds the A3 column per \$15m of all-in cost.

\begin{table}[ht]\centering\small
\caption{AI\_facilities investment column, A3, per MW of IT load}\label{tab:faccol}
\resizebox{\textwidth}{!}{\begin{tabular}{>{\raggedright\arraybackslash}p{4.6cm}c>{\raggedright\arraybackslash}p{7.4cm}}
\toprule
Element & \$m & Producer (NAICS) \\
\midrule
Shell, core, architectural & 1.08 & Construction (BEA 2332A0) \\
Contractor preliminaries and fees & 1.20 & Construction \\
Electrical installation (40\%) & 2.30 & Construction (238210 electrical contractors) \\
Mechanical installation (35\%) & 1.39 & Construction (238220) \\
Site works and civil & 0.60 & Construction \\
Customer substation & 0.30 & Half Construction, half 335311/335313 \\
Electrical equipment (60\%) & 3.46 & See the split below \\
Mechanical equipment (65\%) & 2.57 & 333415 HVAC, chillers, CRAHs, CDUs; 333914 pumps; 333413 fans; some 332 pipe and manifolds \\
Design, engineering, commissioning & 0.90 & 541330, 5413 \\
Contingency & 0.60 & Pro rata \\
Land & 0.60 & Not produced; excluded \\
\midrule
\textbf{Produced total} & \textbf{14.4} & \textbf{Construction 49\%, Electrical 26\%, Machinery 19\%, Prof.\ services 6.5\%} \\
\bottomrule
\end{tabular}}
\end{table}

\paragraph{Air-cooled halls.} For A1, T\&T's air-cooled shares with the same equipment ratios give
Construction 53, Electrical 29, Machinery 12, Prof.\ services 6. With a 70 percent equipment
share, A3 becomes 43/30/20/6.5.

\paragraph{Electrical equipment by maker.} Shares of electrical equipment, A3 (A1 in brackets):
\begin{itemize}
  \item generators, 335312 (engines 333618): 29 (31);
  \item UPS, PDUs and miscellaneous, 335999: 22 (21);
  \item switchgear, switchboards and busway, 335313: 22 (22);
  \item transformers, 335311: 9 (8);
  \item cable and wire, 335929: 7 (7);
  \item batteries, 335911: 6 (5);
  \item controls and other, 335314 and 334513: 5.
\end{itemize}
These replace the unsourced 30/20/15/15/10/5/5 in \texttt{upstream\_concentration.py}, which gave
too much to switchgear, transformers and batteries.

\paragraph{How electrical capex compares by generation.} Installed electrical capex is about
\$5--6m per MW in every generation (A3 \$6.0m, range 4.5--8). The bottom-up build, from
Schneider unit costs indexed by the producer price index (PPI), comes in 15--25 percent below
the consultants. PPIs miss the 2023--26 backlog premiums. Transformer prices rose 76 percent from
2019 to 2025 (PPI; Wood Mackenzie and NIAC agree), and lead times are 120--150 weeks.

\paragraph{Imports of electrical equipment, 2024--25} (share of value imported):

\begin{center}\small
\begin{tabular}{lcp{8cm}}
\toprule
Equipment & Best (range) & Note \\
\midrule
Transformers & 40\% (30--55) & Large power transformers about 75--80\%; distribution transformers 25--50\% \\
Switchgear & 25\% (15--45) & Medium-voltage switchgear 9--15\% \\
Generator sets & 30\% (20--50) & BEA's 64\% looks too high for large diesel sets \\
UPS & 50\% & \\
Lithium-ion batteries & 50\% & \\
\bottomrule
\end{tabular}
\end{center}

\paragraph{How the national accounts record it.} This matters for decision R4:
\begin{itemize}
  \item Census C30 and BEA treat integral electrical and mechanical plant as structures.
  \item BEA added a data-center line to NIPA Table 5.4 in September 2025.
  \item The UK statistics office (ONS, 2026) books racks, PDUs and stand-alone backup generators
    as machinery.
\end{itemize}
The model's Construction sector has the 2017 average recipe, which buys almost no switchgear or
transformers. Routing equipment directly to its producers therefore preserves the supply chain.

\subsection{Facility operations}

\subsubsection{Electricity}

\begin{table}[ht]\centering\small
\caption{Electricity per MW of IT load per year}\label{tab:elec}
\resizebox{\textwidth}{!}{\begin{tabular}{lccccc}
\toprule
 & A0 & A1 & A2 & A3 & A4 \\
\midrule
Operating PUE & 1.70 & 1.15 & 1.25 & 1.15 & 1.10 \\
Utilisation (average draw / provisioned) & 0.50 & 0.65 & 0.70 & 0.75 & 0.80 \\
GWh per MW-IT-year & 7.4 & 6.5 & 7.7 & 7.6 & 7.7 \\
Price, \$/MWh & 110 & 75 & 80 & 85 & 85 \\
\textbf{Cost, \$k per MW-IT-year} & \textbf{819} & \textbf{491} & \textbf{613} & \textbf{642} & \textbf{655} \\
\bottomrule
\end{tabular}}
\end{table}

Energy per MW of IT is roughly flat after the legacy generation, because rising utilisation
offsets falling PUE. Hyperscaler fleet PUEs agree with each other: Google 1.09, Meta 1.08, AWS
1.14--1.15, Microsoft 1.16--1.17. Industry surveys are higher (Uptime 1.36) because they include
colocation and enterprise halls. No measured PUE exists yet for a GB200-class site.

On-site generation is about 90 GW announced but only 2 GW operating. Treat it as a scenario and
as utility-type capital (turbines and engines, 333611 and 333618; gas, 211000 and 486000), not
as part of AI\_facilities.

\subsubsection{Water: old versus new designs}

\begin{table}[ht]\centering\small
\caption{Water per MW of IT load}\label{tab:water}
\resizebox{\textwidth}{!}{\begin{tabular}{lccccc}
\toprule
 & A0 & A1 & A2 & A3 & A4 \\
\midrule
On-site use, L per kWh of IT energy & 2.0 & 0.30 & 0.50 & 0.05 & 0.02 \\
Water and sewer cost, \$k per MW-year & 26 & 4 & 7 & 0.6 & 0.2 \\
Water cost / electricity cost & 3.4\% & 0.7\% & 1.1\% & 0.1\% & $\approx$0 \\
\bottomrule
\end{tabular}}
\end{table}

New designs cut water in two steps:
\begin{itemize}
  \item Outside-air hyperscale cooling cut on-site water 5--10 times at almost no energy cost.
  \item Warm-water liquid cooling with dry coolers takes it to about zero with no energy
    penalty: PUE 1.10, against 1.14 with cooling towers.
\end{itemize}
The saving comes from the heat-rejection choice, not from liquid cooling as such. An A3 hall on
cooling towers still uses about 2 L/kWh.

Legacy halls used towers because going waterless cost about 0.3 points of PUE. That is roughly
\$138k of extra electricity to save \$38k of water per MW-year.

Counting the water used to generate the electricity (about 4.4 L/kWh), total water per unit of
IT falls only 35--40 percent from A0 to A3.

The ``industry average 1.8 L/kWh'' often quoted is a 2016 LBNL assumption per kWh of total site
energy, not a measurement.

\subsubsection{Labour, maintenance and other running costs}

\begin{table}[ht]\centering\small
\caption{Facility running costs other than electricity, \$k per MW of IT load per year}\label{tab:opex}
\resizebox{\textwidth}{!}{\begin{tabular}{>{\raggedright\arraybackslash}p{3.8cm}ccccc>{\raggedright\arraybackslash}p{3.4cm}}
\toprule
 & A0 & A1 & A2 & A3 & A4 & Supplier \\
\midrule
Facility staff, FTE per MW & 5 & 1.3 per hall; 0.6 per campus & 0.7 & 0.5 & 0.3 & Half own staff; half 5612/5616/5617 \\
Own payroll & 400 & 133 & 90 & 60 & 40 & Labour \\
Plant maintenance & 150 & 70 & 100 & 120 & 150 & 230301 repairs; some 811300 \\
Security & 150 & 43 & 30 & 20 & 15 & 5616 \\
Cleaning & 40 & 13 & 7 & 5 & 3 & 5617 \\
Insurance & 10 & 10 & 12 & 15 & 20 & 524 \\
Property tax, after abatements & 80 & 50 & 70 & 80 & 95 & Tax on production \\
Software, IT and network operations & 100 & 60 & 50 & 50 & 55 & 517, 518, 5415 \\
Other (generator fuel and testing, consumables) & 40 & 25 & 25 & 30 & 30 & 324110, misc. \\
\midrule
\textbf{Total} & \textbf{$\sim$1{,}000} & \textbf{$\sim$410} & \textbf{$\sim$390} & \textbf{$\sim$380} & \textbf{$\sim$410} & \\
\bottomrule
\end{tabular}}
\end{table}

Staffing agrees across sources at about 1.3 FTE per MW for a 30--40 MW hall (Virginia JLARC,
KPMG, Uptime 2007). At gigawatt scale the sources diverge, from Meta at 0.2 to operator
announcements at 0.9--1.5. BLS employment in NAICS 518210 is not a measure of facility staff.
Most facility staff work for the parent company or for facilities-support and security firms.

\subsubsection{Service lives}

\begin{itemize}
  \item Shell: about 40 years (BEA offices: 36 years, $\delta=0.0247$).
  \item Electrical switchgear, transformers and busway: about 30 years.
  \item UPS electronics: about 10 years. Lithium-ion batteries: 12. VRLA batteries: 4, which is
    a replacement flow in legacy halls.
  \item Mechanical plant: 18--20 years.
  \item Blended facility $\delta$: 0.078 with plant at 0.088 (BEA's 1.65/L rule with a 19-year life); 0.085--0.09 with plant at 0.10.
\end{itemize}

\subsection{IT capital}

\begin{table}[ht]\centering\small
\caption{IT capital per generation}\label{tab:it}
\resizebox{\textwidth}{!}{\begin{tabular}{lccccc}
\toprule
 & A0/A1 & A2 H100 & A3 GB200 & A3 GB300 & A4 Rubin \\
\midrule
Rack or server price & \$7.5--10k per server & \$190--270k per server & \$3.1m & \$4.5m & $\sim$\$7m \\
All-in per GPU & -- & \$35k & \$54k & \$76k & $\sim$\$115k \\
\textbf{\$m per MW of IT} & \textbf{15} (10--22) & \textbf{24} (21--28) & \textbf{27} (26--30) & \textbf{37} (31--42) & \textbf{$\sim$50} (35--65) \\
Network share of IT capex & 15\% & 15\% & 14\% & 13\% & 12\% \\
Storage share & 15\% & 4\% & 4\% & 4\% & 5--15\% \\
Integration share & 3\% & 3\% & 3\% & 3\% & 2--4\% \\
NVIDIA gross profit / IT capex & -- & $\sim$55\% & 55--60\% & 55--60\% & 55--60\% \\
Geometric $\delta$ & 0.30 & 0.27 & 0.27 & 0.27 & 0.27--0.33 \\
Import share & \multicolumn{5}{c}{0.85--0.95 (servers); networking 0.80--0.95; storage 0.60--0.80} \\
\bottomrule
\end{tabular}}
\end{table}

Confidence is highest for H100 and GB200, which have several primary analyst sources, vendor
specifications and filings. It is medium for GB300 and low for A0/A1 per MW and for A4.

Cost per rack rose about 15 times from CPU racks to GB200, but cost per MW only about 2 times,
because rack power rose about 12 times.

\paragraph{Composition, A3.} Shares of cluster IT capex:
\begin{itemize}
  \item GPU modules including HBM: 55;
  \item the rest of the compute tray: 14;
  \item NVLink switches and rack infrastructure: 10;
  \item scale-out network: 14;
  \item storage: 4;
  \item integration: 3.
\end{itemize}
NVIDIA content is over 92 percent of the rack, at a 71--75 percent gross margin. With a closed
economy, that margin becomes value added of Computers. This is why the import treatment matters.

\paragraph{Investment column of AI\_compute, at purchasers' value.} Central values (range):
\begin{itemize}
  \item computers (334111, 334112): 0.77;
  \item communications equipment (334210; networking): 0.13;
  \item integration, computer systems design (541512): 0.03 (0.02--0.05);
  \item wholesale: 0.06 (0.03--0.10);
  \item transport: 0.01;
  \item retail: 0.
\end{itemize}
This replaces the average-buyer margins in the BEA 2017 bridge, where wholesale is 0.19. US
server imports reached \$164bn in 2025, 91 percent from Mexico and Taiwan. 58 percent of server
exports are re-exports: GPUs and boards go out and come back as servers.

\paragraph{Depreciation.} Five independent kinds of evidence give $\delta\approx0.20$--$0.30$ for
GPU servers:
\begin{itemize}
  \item \textbf{Accounting lives} of 5--6 years in 31 SEC filings. Amazon cut part of its fleet
    to 5 years in 2025, citing AI. Neoclouds use 4--6 years.
  \item \textbf{Rents across chip generations:} about 0.27 log points per year of vintage.
  \item \textbf{Resale prices by age:} 0.18--0.28.
  \item \textbf{Price-performance:} 0.28--0.40 log per year (Epoch).
  \item \textbf{Physical failure is small:} replacements of 0.4--0.8 percent a year. The widely
    quoted 9 percent counts job interruptions, not dead GPUs.
\end{itemize}
The 2023--25 fall in H100 rents (0.35--0.48 a year) includes the unwinding of a scarcity premium
and should not be read as $\delta$.

\paragraph{Compute's operating inputs other than power and space.} Per \$1 of GPU-cloud revenue:
\begin{itemize}
  \item software: 0.02;
  \item telecom and transit: 0.01;
  \item maintenance and spares: 0.015;
  \item purchased sales, marketing and administrative services: 0.025;
  \item labour, excluding the labs: 0.06--0.12;
  \item property tax: about 0.01, a production tax.
\end{itemize}
Sources: CoreWeave, Nebius and IREN filings. The Bank of Italy's network, maintenance and
``other'' lines are 2--4 percent of CoreWeave's realised revenue per rack, which is consistent.
Its list rental prices are 4--6 times realised revenue.

\subsection{Per \$1 of AI services}\label{sec:perdollar}

The recipe is scaled from per MW of IT load to \$1 of AI services in three steps:
\begin{enumerate}
  \item The facility's output (powered space) is its capital at user cost $r+\delta_f$, plus
    electricity, other purchases, labour and property tax.
  \item Compute costs its IT capital at user cost $r+\delta_z$, plus space, plus operating inputs
    (purchases 6.5, staff 9 and property tax 1 percent of its cost).
  \item Korinek and McKelvey's closure turns compute into AI services: services revenue is 1.5
    times inference compute, the labs' labour is 0.1 times inference, and the margin pays for
    training compute. Compute is therefore 0.933 of each \$1.
\end{enumerate}
Table~\ref{tab:rec:perdollar} reports the result by generation, with the 2026 rack recipe for
comparison.

\begin{table}[ht]\centering\small
\caption{Recipe per \$1 of AI services, by generation}\label{tab:rec:perdollar}
\resizebox{\textwidth}{!}{%
\begin{tabular}{lccccccc}
\toprule
 & A0 & A1 & A2 & A3 & A3 GB300 & A4 & 2026 rack \\
\midrule
$z$: labour (labs + staff running the computers) & 0.151 & 0.151 & 0.151 & 0.151 & 0.151 & 0.151 & 0.067 \\
$z$: IT capital services & 0.487 & 0.537 & 0.588 & 0.589 & 0.632 & 0.650 & 0.669 \\
$z$: space from $f$ & 0.297 & 0.247 & 0.196 & 0.195 & 0.152 & 0.134 & 0.264 \\
$z$: other purchases & 0.065 & 0.065 & 0.065 & 0.065 & 0.065 & 0.065 & 0 \\
$f$: electricity & 0.080 & 0.053 & 0.049 & 0.046 & 0.036 & 0.027 & 0.080 \\
$f$: other purchases & 0.050 & 0.024 & 0.018 & 0.017 & 0.013 & 0.011 & 0.032 \\
$f$: labour & 0.039 & 0.014 & 0.007 & 0.004 & 0.003 & 0.002 & 0.032 \\
$f$: capital services & 0.128 & 0.155 & 0.122 & 0.128 & 0.100 & 0.093 & 0.120 \\
\addlinespace
IT capital per \$1 a year & 1.49 & 1.65 & 1.94 & 1.94 & 2.08 & 2.13 & 2.29 \\
Facility capital per \$1 a year & 1.12 & 1.31 & 1.05 & 1.07 & 0.83 & 0.76 & 1.03 \\
$\delta_{IT}$ / $\delta_f$ & .285/.072 & .285/.076 & .261/.074 & .261/.078 & .262/.078 & .264/.081 & .25/.075 \\
Electricity share of $f$'s cost & 0.27 & 0.21 & 0.25 & 0.23 & 0.23 & 0.21 & 0.30 \\
AI services per MW-IT-year, \$m & 10.3 & 9.3 & 12.5 & 14.1 & 18.1 & 23.8 & -- \\
\addlinespace
$f$ investment: Constr./Elec./Mach./Prof. & 55/25/13/7 & 52/29/11/7 & 51/28/13/7 & 49/26/17/7 & 49/26/17/7 & 48/27/17/7 & 33/40/27/0 \\
$z$ investment: Computers/integration/wholesale & \multicolumn{6}{c}{90 / 3 / 6 (transport 1)} & 78/0/22 \\
\bottomrule
\end{tabular}}
\end{table}

Newer generations spend a larger share of each dollar on IT capital and a smaller share on
space. Electricity falls from 8 cents per dollar (legacy) to under 3 (next generation), because
chips get more expensive per watt faster than power gets cheaper.

Two decisions are kept unchanged for the authors (\texttt{recipe\_plan.md}, R1--R2): the labs'
labour (Korinek and McKelvey's 0.1 of inference, which the evidence puts about 10 times too low)
and the treatment of training as a current cost.

One cross-check is open. At user cost on new capital, compute costs about \$13m per MW-year.
Epoch's total cost of ownership and CoreWeave's revenue per MW are about \$8.5m.

\subsection{Cross-checks}

\begin{longtable}{>{\raggedright\arraybackslash}p{3.2cm}>{\raggedright\arraybackslash}p{5.6cm}>{\raggedright\arraybackslash}p{1.6cm}>{\raggedright\arraybackslash}p{4.0cm}}
\caption{Where independent sources agree and disagree}\label{tab:xcheck}\\
\toprule
Number & Sources & Agree? & Why they differ / verdict \\
\midrule\endfirsthead
\toprule
Number & Sources & Agree? & Verdict \\
\midrule\endhead
IT share of campus capex & Van Nieuwerburgh 68; Epoch 69; Fed 70 (IT-inclusive); Microsoft PP\&E 57 servers; CNT rack + network 72 & Yes & 68--72 \\
Facility hard cost per MW & T\&T; JLL; KPMG; Schneider & Yes, once scope and year match & \$12m (A3) \\
Facility all-in cost per MW, A3 & Landlords 8--13; DLR 14.1; C\&W 17.6 & Spread & Scope: brownfield vs land, soft costs and contingency \\
Shell share & 9--35\% across sources & No & Some include site work; use 9--14\% plus 3--6\% site \\
Equipment vs installation & Schneider WP61 (only primary); KPMG; Eaton content & Single source & Gap narrowed, not closed \\
Electrical capex per MW & Bottom-up PPI build; T\&T $\times$ share; C\&W & Partly & PPIs miss backlog premiums; lean to the consultants \\
Eaton AI content & \$2.4m/MW vs bottom-up \$1.3m & No & Possibly IT-side power; unresolved \\
Transformer price inflation & PPI +76\%; Wood Mackenzie +77--95\%; NIAC +80\% & Yes & \\
Transformer import share & Own calculation 26.5\% (2017), $\ge$38.5\% (2024); BEA 29\%; BIS 82\% of large units & Yes by size class & \\
Generator import share & Own calculation $\le$44\%; BEA 64\% & No & BEA too high for large sets \\
Hyperscale PUE & Google, Meta, AWS, Microsoft 1.08--1.17; LBNL 1.22; Uptime 1.36 & Operators yes & Surveys include colocation and enterprise \\
Hyperscale water (L/kWh) & LBNL 0.32; Microsoft 0.27--0.30; Meta 0.20; Virginia utilities 0.2--0.3 & Yes & Withdrawal vs consumption explains a 1.3$\times$ gap \\
Staff per MW, 30 MW hall & JLARC; KPMG; Uptime & Yes & 1.3 \\
Staff per MW, 1 GW campus & Meta 0.2; Crusoe 0.47; Oracle and Microsoft announcements 0.9--1.5 & No & 0.5 (0.3--1.0) \\
Facility running costs other than electricity & KPMG; ITK; Epoch; Digital Realty 0.33--0.41m; CNT 0.64m & Yes except CNT & CNT includes network (belongs to $z$) and an unsourced line \\
Electricity share of facility cost & Own A3 25\%; engineering notes 25--30\%; Digital Realty 23\% of revenue & Yes & \\
GB200 rack price & CNT \$3.0m; SemiAnalysis \$3.1m & Yes & Independent \\
Non-rack share of IT & SemiAnalysis (2 articles); Cottier; Epoch & Yes & Van Nieuwerburgh's 12.5\% is low and unsourced \\
GB300 price & Van Nieuwerburgh \$4.0m; purchase orders \$5.0m; press \$6--6.5m & No & $\pm$25\% \\
IT $\delta$ & SEC lives; rents; resale; Epoch & Yes & 0.26 (0.20--0.32) \\
Wholesale margin on servers & BEA PEQ 19\%; distributors 7--11\%; direct purchase & No & 6\% for hyperscalers \\
\bottomrule
\end{longtable}

\paragraph{Mis-quotes found and corrected.}
\begin{itemize}
  \item ``1.8 L/kWh industry average'': a 2016 LBNL assumption, not a measurement.
  \item ``9 percent annual GPU failure rate'': a press calculation on Meta's interruption table.
    It counts interruptions, and the paper itself has a typo in the table.
  \item ``GPUs last 1--3 years'': an anonymous social-media post.
  \item Van Nieuwerburgh applies a per-IT-MW cost to total facility power.
  \item Epoch's facility cost re-quotes Turner \& Townsend.
  \item IEA's ``84 percent training load factor'' is not in the LBNL report it cites.
\end{itemize}

\subsection{Imports of AI capital: IT equipment, facility gear and the grid}\label{sec:rec:imports}

\paragraph{The current treatment.} The baseline economy is closed, following the vom
Lehn--Winberry convention. The paper's imports alternative opens one flow only:
\begin{itemize}
  \item Imports supply 95 percent of what AI\_compute buys from Computers, the BEA 2017 import
    share of computers.
  \item A linear technology turns an export basket with the composition of 2017 consumption into
    imported equipment at a fixed rate.
  \item Trade balances every year.
\end{itemize}
Facility gear (electrical and mechanical equipment) and the grid equipment that Utilities buys
stay domestic. The paper's conclusion lists imports of data-center electrical equipment as a next
step. This section sets out how to add them. Nothing here is implemented yet.

\paragraph{How big the flows are.} In the steady state of the A3 recipe, per \$1 of AI services:
\begin{itemize}
  \item AI\_facilities invests $\delta_fK_f\approx0.080$ a year: 0.021 in electrical equipment
    and 0.013 in machinery.
  \item AI\_compute invests $\delta_zK_z\approx0.50$ a year, of which 0.45 is Computers
    (servers, networking and storage).
\end{itemize}
At the import shares below, imported facility gear is about 0.010 per \$1 of AI services, about
1 percent of AI revenue. Imported IT is about 0.40. Gear imports therefore barely move GDP,
consumption or AI investment.

They matter for industries. Electrical equipment is one of the paper's headline winners, and in
the IT-imports alternative it is the industry that takes the place of Computers. If about 40
percent of the AI-driven demand for electrical gear goes abroad, that gain shrinks accordingly.

\begin{table}[ht]\centering\small
\caption{Import shares of AI capital (share of value imported)}\label{tab:rec:imports}
\resizebox{\textwidth}{!}{%
\begin{tabular}{llccl}
\toprule
Item & NAICS & BEA 2017 & 2024--25 (best, range) & Note \\
\midrule
\multicolumn{5}{l}{\textit{AI\_compute: IT equipment}} \\
Servers & 334111 & 0.95 & 0.90 (0.85--0.95) & Mexico and Taiwan are 91\% of 2025 server imports \\
Networking & 334210 & -- & 0.875 (0.80--0.95) & Vietnam, Thailand \\
Storage & 334112 & -- & 0.70 (0.60--0.80) & Thailand \\
\addlinespace
\multicolumn{5}{l}{\textit{AI\_facilities: electrical equipment (share of it in brackets)}} \\
Generator sets (29\%) & 335312 & 0.64 & 0.30 (0.20--0.50) & BEA too high for large diesel sets \\
UPS, PDUs (22\%) & 335999 & 0.55 & 0.50 (0.35--0.60) & Proxy: static converters \\
Switchgear, busway (22\%) & 335313 & 0.26 & 0.25 (0.15--0.45) & Medium voltage 9--15\% \\
Transformers (9\%) & 335311 & 0.29 & 0.40 (0.30--0.55) & $\ge$38.5\% in 2024; large units 75--82\% \\
Cable and wire (7\%) & 335929 & 0.91 & not re-checked & BEA 2017 value \\
Batteries (6\%) & 335911 & $\approx$0.36 & 0.50 (0.40--0.70) & Li-ion; VRLA $\approx$0.35 \\
Controls (5\%) & 335314 & -- & $\approx$0.45 & Broad switchgear and controls, 43--50\% \\
\textbf{Weighted electrical} & & $\approx$\textbf{0.50} & $\approx$\textbf{0.40} & \\
\addlinespace
\multicolumn{5}{l}{\textit{AI\_facilities: mechanical equipment}} \\
HVAC, chillers, CDUs & 333415 & 0.14 & $\approx$0.15 & Pumps and fans not studied \\
\addlinespace
\multicolumn{5}{l}{\textit{Utilities: grid equipment}} \\
Large power transformers & 335311 & -- & 0.75--0.82 & BIS 2020; Wood Mackenzie \\
Distribution transformers & 335311 & -- & 0.20--0.50 & BIS 18--21\% (2019) vs Wood Mackenzie 50\% \\
Medium- and high-voltage switchgear & 335313 & 0.26 & 0.15 & \\
\bottomrule
\end{tabular}}
\end{table}

\paragraph{How to implement it in the model.} Generalise the single import weight to one weight
$m_{ij}$ for each pair of supplying sector $i$ and investing sector $j$. The price of $j$'s
investment bundle becomes
\[
\tilde P^k_{j}=\Big(\sum_i (1-m_{ij})\,G^I_{ij}\,P_i^{1-\sigma_I}+\Big(\sum_i m_{ij}G^I_{ij}\Big)P_{imp}^{1-\sigma_I}\Big)^{1/(1-\sigma_I)},
\]
where $P_{imp}$ is the price of the export basket that pays for imports. The rest follows
directly:
\begin{itemize}
  \item Imports by sector $j$ are
    $I^{imp}_j=\big(\sum_im_{ij}G^I_{ij}\big)(P_{imp}/\tilde P^k_j)^{-\sigma_I}I_j$.
  \item Exports are the consumption-weighted basket that pays for $\sum_jI^{imp}_j$.
  \item The resource constraint and GDP add up over $j$.
\end{itemize}
In the code, the scalar import weight becomes a matrix, filled from import shares stored by NAICS
code in \texttt{ai\_recipe.py} and aggregated with the recipe. Two choices remain:
\begin{enumerate}
  \item \textbf{Vintage of the shares.} BEA 2017 detail is consistent with the vintage of the
    current IT exercise. The 2024--25 trade data match a 2025--26 recipe. The weighted
    electrical share is about 0.50 on BEA 2017 and 0.40 on 2024--25; the main difference is
    generator sets. Using 2024--25 for both IT and gear moves the IT share only from 0.95 to about
    0.90.
  \item \textbf{Utilities.} The grid bottleneck runs through the equipment Utilities buys to
    expand transmission and distribution: transformers and switchgear, 15 percent of its 2017
    investment column. Importing them means opening an existing sector's investment column, which
    departs from the closed-economy convention for the 37 sectors. There are two ways to do it:
    \begin{itemize}
      \item \emph{Average shares.} Utilities imports its electrical equipment at all dates,
        including 2017. Then the 2017 steady state has trade, and the export basket must be
        calibrated to it.
      \item \emph{Marginal shares.} Only purchases above the 2017 level are imported. This keeps
        the 2017 calibration unchanged but is less standard.
    \end{itemize}
    The first is cleaner and makes the grid, and its import dependence, part of the exercise.
    The authors plan to include Utilities if the extension proceeds.
\end{enumerate}

\paragraph{What it would add to the paper.}
\begin{itemize}
  \item A version of the imports alternative with IT, facility gear and grid equipment.
  \item How much of the Electrical sector's gain leaks abroad.
  \item Whether the ranking of winners survives.
  \item A natural tariff scenario: a tax on imported AI capital. See the Quintana et al.\ (2026)
    paper on tariffs as taxes on capital in \texttt{PandP/}.
\end{itemize}
Run time is about 10 minutes per scenario.

\subsection{Using the recipe with the 2017 detailed tables}

The recipe is specified at 6-digit NAICS wherever the sources allow. The proposal is to store it
as a vector over the 2017 detailed commodities:
\begin{itemize}
  \item one vector per generation and per sector (running costs and investment column);
  \item with the source, year and range attached to each cell;
  \item and let the loader aggregate it to the 37 sectors.
\end{itemize}
The same data then serve the paper as it stands, a detailed or split version of the model (the
\texttt{Construction\_AI}-type splits already in \texttt{model.py}), and import shares by
commodity, which differ a lot within Electrical.

Two caveats:
\begin{itemize}
  \item Several NAICS codes map to aggregated BEA detail codes, e.g.\ pumps and batteries. The
    NAICS-to-BEA mapping is still to be done.
  \item The 2017 tables have no data-center industry. Hosting (518200) is mostly payment
    processing, and data-center buildings sit in 2332A0. The detailed tables supply the
    suppliers, not a data-center column to copy.
\end{itemize}

\subsection{Open gaps}
\begin{itemize}
  \item Equipment versus installation rests on one primary source (Schneider WP61).
  \item No measured PUE or water use for GB200-class sites. No cost data on liquid-cooling
    maintenance.
  \item The split of mechanical equipment among HVAC, pumps and fans is not quantified.
  \item Insurance rates; the size of owner-furnished equipment beyond KPMG; grid-upgrade cost per
    kW.
  \item Prices for GB300 and Rubin rest on press reports and one blog. SemiAnalysis item-level
    tables are paywalled.
  \item The share of hyperscaler servers bought directly from manufacturers.
  \item The labs' cost structure (decisions R1--R2).
\end{itemize}

\subsection{Main sources}
\begin{itemize}
  \item \textbf{Cost consultants:} Turner \& Townsend, Data Centre Construction Cost Index
    (2019--25); KPMG (2026), benchmarking capex and opex; JLL 2026 outlooks; Cushman \& Wakefield
    (partly verified).
  \item \textbf{Engineering and vendors:} Schneider Electric white papers WP61, WP75, WP218,
    WP229 and WP282; Uptime Institute (Koomey et al.\ 2008; annual surveys); Eaton; NVIDIA DGX and
    NVL72 specifications and its 800 V DC paper.
  \item \textbf{National labs, agencies and standards:} LBNL (Shehabi et al.\ 2024); EPRI; EIA;
    The Green Grid WP35; ASHRAE; BIS and DOE on transformers; Mytton (2021); Li et al.\ (2023).
  \item \textbf{Official and filings:} Virginia JLARC (2024); BEA depreciation rates and NIPA 5.4;
    ONS (2026); UN Comtrade and Census trade data; SEC filings of Amazon, Microsoft, Alphabet,
    Meta, Oracle, CoreWeave, Nebius, IREN, WhiteFiber, Nscale, Digital Realty and Equinix.
  \item \textbf{Analysts and measurement:} SemiAnalysis (GB200 architecture, AI cloud total cost
    of ownership (TCO), GPU index); Epoch AI (hardware trends, 1 GW TCO); Silicon Data and Ornn
    rent indices; Meta, NCSA, BigScience and ByteDance reliability studies; Cottier et al.\
    (2025).
\end{itemize}
Files:
\begin{itemize}
  \item Reports and observations: \texttt{notes/recipe\_research/*.md} and \texttt{*\_obs.csv}.
  \item PDFs: \texttt{References/engineering/}, in four subfolders (facility capex, power and
    energy, operations, IT equipment).
\end{itemize}
